\kbvsechead{ವಿಶೇಷ ರೀತಿಯಾಗಿ ಸ್ತ್ರೀಲಿಂಗ ರೂಪ ಶಬ್ದಗಳ ಉತ್ಪತ್ತಿ:}

\indent ಅಪ್ಪ \ensuremath{\times} ಅಮ್ಮ, ತಂದೆ \ensuremath{\times} ತಾಯಿ, ಚಿಕ್ಕಪ್ಪ \ensuremath{\times} ಚಿಕ್ಕಮ್ಮ, ತಮ್ಮ \ensuremath{\times} ತಂಗಿ, ವಿಧುರ \ensuremath{\times} ವಿಧವೆ, ರಾಜ \ensuremath{\times} ರಾಣಿ -- ಇವೇ ಮೊದಲಾದವುಗಳು ಸಾಮಾನ್ಯ ರೀತಿಯ ಪುಲ್ಲಿಂಗ-ಸ್ತ್ರೀಲಿಂಗ ಶಬ್ದಗಳು. ಆದರೆ ಕೆಲವು ವಿಶೇಷ ಸಂದರ್ಭಗಳಲ್ಲಿ ಪುಲ್ಲಿಂಗ ನಾಮಪ್ರಕೃತಿಗಳ ಮೇಲೆ ಎ, ಇ, ಅಳು, ಇತ್ತಿ , ಗಿತ್ತಿ, ಇತ್ತಿ ಮೊದಲಾದ ಪ್ರತ್ಯಯಗಳು ಸೇರಿ ಸ್ತ್ರೀಲಿಂಗ ರೂಪ ಶಬ್ದಗಳು ಉತ್ಪತ್ತಿಯಾಗುವುವು.\par

\vspace{3pt plus 2pt}\noindent\hspace*{1.2em}\kbvlabel{ಉದಾ}\par\vspace{1pt}

\begin{kbvderiv}
\vspace{4pt plus 3pt minus 1pt}
\begin{center}\begin{adjustbox}{max width=\textwidth}
{\renewcommand{\arraystretch}{1.22}\setlength{\tabcolsep}{5pt}
\begin{tabular}{l c l l}
ಗುಣವಂತ & $\times$ & ಗುಣವಂತೆ & \ldots\ (ಗುಣವಂತ + ಎ) \\
ಶಿಕ್ಷಕ & $\times$ & ಶಿಕ್ಷಕಿ & \ldots\ (ಶಿಕ್ಷಕ + ಇ) \\
ಮಗ & $\times$ & ಮಗಳು & \ldots\ (ಮಗ + ಅಳು) \\
ಬೀಯಗ & $\times$ & ಬೀಯಗಿತಿ & \ldots\ (ಬೀಯಗ + ಇತಿ) \\
ಕುಂಬಾರ & $\times$ & ಕುಂಬಾರಗಿತ್ತಿ & \ldots\ (ಕುಂಬಾರ + ಗಿತ್ತಿ) \\
ಹೂವಾಡಿಗ & $\times$ & ಹೂವಾಡಿಗಿತ್ತಿ & \ldots\ (ಹೂವಾಡಿಗ + ಇತ್ತಿ) \\
\end{tabular}}
\end{adjustbox}\end{center}
\vspace{4pt plus 3pt minus 1pt}
\end{kbvderiv}

\kbvexercise{ಅಭ್ಯಾಸ ಪತ್ರಿಕೆ ೯}

\vspace{2.5pt plus 3pt minus 0.5pt}\noindent\hspace*{1.2em}\kbvnum{(1)}\hspace{0.55em}ಲಿಂಗಗಳಲ್ಲಿ ಮುಖ್ಯವಾಗಿ ಎಷ್ಟು ಪ್ರಭೇದಗಳಿವೆ? ಅವು ಯಾವುವು? ಉದಾಹರಣೆಗಳ ಸಹಿತ ತಿಳಿಸಿ.\par

\vspace{2.5pt plus 3pt minus 0.5pt}\noindent\hspace*{1.2em}\kbvnum{(2)}\hspace{0.55em}ಈ ಪದಗಳ ಪುಲ್ಲಿಂಗ ರೂಪಗಳನ್ನು ತಿಳಿಸಿ:\par

\indent ಭಿಕ್ಷುಕಿ, ಸಂನ್ಯಾಸಿನಿ, ಗೆಳತಿ, ಕುಂಟಿ, ಯುವತಿ, ಕಿವುಡ, ಮಾಟಗಾತಿ, ಮೂಗಿ, ಹೆಳವ, ಕನ್ಯೆ.\par

\vspace{2.5pt plus 3pt minus 0.5pt}\noindent\hspace*{1.2em}\kbvnum{(3)}\hspace{0.55em}ಈ ಪದಗಳ ಸ್ತ್ರೀಲಿಂಗ ರೂಪಗಳನ್ನು ತಿಳಿಸಿ:\par

\indent ಗಾಣಿಗ, ಸೇವಕ, ದಾಸ, ಸೊಗಸುಗಾರ, ಪೂಜಾರಿ, ಮುದುಕ, ಹಣವಂತ, ಅಯ್ಯ, ಸಹೋದರ, ವರ.\par

\vspace{2.5pt plus 3pt minus 0.5pt}\noindent\hspace*{1.2em}\kbvnum{(4)}\hspace{0.55em}ಈ ಪದಗಳು ಯಾವ ಲಿಂಗಕ್ಕೆ ಸೇರಿದವುಗಳು?\par

\indent ಗಂಡುಮಗು, ದೇವರು, ಪೂಜಾರಿ, ಸಲಗ, ಅಮ್ಮ, ಗುಮಾಸ್ತೆ, ಚಂದ್ರ, ಗುಲಾಬಿ, ಹೋರಿಗಳು.\par
